
\chapter{کارهای پیشین}
\label{ch:literature}

این فصل پیشینه پژوهشی مرتبط را در چهار محور بررسی می‌کند: سامانه‌های
گفت‌وگوی گفتاری آبشاری در برابر سرتاسری، مدل‌های گفتار‌به‌گفتار متن‌باز و
قابلیت‌های تمام‌دوطرفه آن‌ها، منابع گفتاری زبان فارسی، و پیشینه تشخیص وقفه و
نوبت‌گیری. در پایان فصل، جایگاه این پروژه در برابر کارهای پیشین مشخص می‌شود.

\section{از زنجیره آبشاری به مدل سرتاسری}
\label{sec:cascade-vs-e2e}

معماری متعارف سامانه‌های گفت‌وگوی گفتاری، زنجیره‌ای از چهار مؤلفه مستقل
است: تشخیص فعالیت گفتاری، بازشناسی گفتار، یک موتور گفت‌وگوی متنی و تبدیل
متن به گفتار. مزیت این معماری، ماژولار بودن آن است: هر مؤلفه را می‌توان
جداگانه بهبود داد، اشکال‌زدایی کرد و با بهترین مدل موجود برای آن زبان
جانشین ساخت. برای زبان‌های کم‌منبع، این ماژولار بودن مزیت تعیین‌کننده‌ای است،
زیرا برای هر مؤلفه احتمال یافتن یک مدل فارسی قابل قبول بیشتر از یافتن یک
مدل گفتار‌به‌گفتار فارسی است.

در برابر آن، نویسندگان \en{Moshi}~\cite{defossez2024moshi} سه ایراد بنیادی
برای معماری آبشاری برمی‌شمارند. نخست، انباشت تأخیر در زنجیره که فاصله میان
دو نوبت گفت‌وگو را به چند ثانیه می‌رساند. دوم، از‌دست‌رفتن اطلاعات
غیرزبانی: چون متن، تنها واسطه میان ورودی و خروجی است، عاطفه، لحن و
صداهای غیرگفتاری که معنا را تغییر می‌دهند از میان می‌روند. سوم، وابستگی به
بخش‌بندی صریح گفتار به نوبت‌ها، که هم‌پوشانی و وقفه و اظهارهای کوتاه تأییدی
را نمی‌تواند نمایش دهد.

مشاهدات فصل~\ref{ch:results} هر دو سوی این استدلال را پشتیبانی می‌کنند. ایراد تأخیر در
اندازه‌گیری‌های مسیر آبشاری دیده می‌شود: زمان تولید کامل پاسخ در زنجیره آبشاری، حتی روی
سخت‌افزار پرتوان، در حدود چند ثانیه است و برای ورودی‌های بلند رشد می‌کند
(بخش~\ref{sec:res-cascade}). در همان حال، مزیت ماژولار بودن نیز دیده
می‌شود: تنها همین معماری بود که در بازه زمانی پروژه به یک سامانه فارسی
کارکردی و قابل ارزیابی رسید.

\section{مدل‌های گفتار‌به‌گفتار متن‌باز}
\label{sec:openmodels}

\subsection{خانواده \en{Mini-Omni}}

مدل \en{Mini-Omni2}~\cite{xie2024miniomni2}، که در تعریف پروژه به‌عنوان نمونه
مدل پایه نام برده شده است، یک دستیار چند‌وجهی بر پایه \en{Qwen2-0.5B} است که
کدگذارهای \en{Whisper} و \en{CLIP} را برای ورودی صوتی و تصویری به کار
می‌گیرد و با راهبرد «تولید موازی با تأخیر متن‌راهنما» گفتار تولید می‌کند.
برای مسئله ما دو ویژگی آن اهمیت دارد: نخست آن‌که سازوکار وقفه در آن
«فرمان‌محور» است و نه یک تصمیم آوایی پیوسته؛ دوم آن‌که این پروژه
آموزش‌دهنده کامل و رسمی گفتار‌به‌گفتار منتشر نکرده است، و بنابراین تطبیق آن
با فارسی مستلزم بازنویسی زنجیره آموزش از صفر بود.

\subsection{خانواده \en{LLaMA-Omni}}

نسخه نخست این خانواده~\cite{fang2025llamaomni} از یک کدگشای گفتاری
غیرخودرگرسیو برای کاهش تأخیر استفاده می‌کرد.
\en{LLaMA-Omni2}~\cite{fang2025llamaomni2} مجموعه‌ای از مدل‌های گفتاری ماژولار
از \N{0.5} تا \N{14} میلیارد پارامتر است که بر پایه مدل‌های
\en{Qwen2.5-Instruct} ساخته شده‌اند، کدگذار \en{Whisper} را برای ورودی و یک
کدگشای گفتاری جویباری خودرگرسیو الهام‌گرفته از
\en{CosyVoice 2}~\cite{du2024cosyvoice2} را برای خروجی به کار می‌برند. این
مدل با تنها \N{200} هزار نمونه گفت‌وگوی چند‌نوبتی آموزش دیده و از
\en{GLM-4-Voice}~\cite{zeng2024glm4voice} که بر میلیون‌ها ساعت گفتار آموزش
دیده بود پیشی گرفته است.

این مدل در آغاز پروژه گزینه اصلی بود، زیرا اندازه \N{0.5} میلیارد
پارامتری آن در بودجه حافظه هدف جای می‌گرفت و پشته آن (کدگذار \en{Whisper} به
اضافه \en{Qwen2.5}) با ابزارهای در دسترس هم‌خوان بود. اما مخزن منتشرشده
تنها استنتاج را پوشش می‌دهد. تلاش نخست این پروژه برای ساخت یک زنجیره آموزش
از این معماری، همان‌طور که در بخش~\ref{sec:early-attempt} شرح داده می‌شود، به
مصنوعی رسید که در واقع بازسازی ورودی را می‌آموخت و نه تولید پاسخ گفت‌وگویی؛
این تجربه، انگیزه اصلی حرکت به‌سوی یک آموزش‌دهنده رسمی شد.

\subsection{\en{Freeze-Omni} و رویکرد مدل منجمد}

\en{Freeze-Omni}~\cite{wang2025freezeomni} راهبرد متفاوتی دارد: وزن‌های مدل
زبانی متنی در سراسر آموزش «منجمد» می‌مانند و تنها کدگذار گفتاری ورودی و
کدگشای گفتاری خروجی آموزش می‌بینند. این طراحی دو مزیت مهم دارد که برای
شرایط کم‌منبع بسیار مرتبط است: هوشمندی مدل پایه دست‌نخورده می‌ماند و مسئله
«فراموشی فاجعه‌بار» پیش نمی‌آید، و آموزش با داده و منابع محاسباتی
بسیار کمتری امکان‌پذیر می‌شود. قابلیت تمام‌دوطرفه نیز در آن با افزودن یک
لایه طبقه‌بندی حالت روی آخرین لایه مدل زبانی به دست می‌آید که پیش‌بینی
می‌کند آیا کاربر گفت‌وگو را قطع کرده است یا نه.

از دید پس‌نگرانه، این رویکرد یکی از امیدبخش‌ترین گزینه‌ها برای فارسی است و
در فصل~\ref{ch:conclusions} به‌عنوان کار آینده پیشنهاد شده است؛ در زمان
تصمیم‌گیری این پروژه، آموزش‌دهنده رسمی و در دسترس آن برای یک زبان جدید
فراهم نبود.

\subsection{\en{Qwen2.5-Omni} و معماری «اندیشمند و گوینده»}

\en{Qwen2.5-Omni}~\cite{xu2025qwen25omni} یک مدل سرتاسری چند‌وجهی است که با
معماری «اندیشمند و گوینده» تولید هم‌زمان متن و گفتار را ممکن می‌سازد:
بخش اندیشمند یک مدل زبانی است که متن تولید می‌کند و بخش گوینده یک مدل
خودرگرسیو دو‌مسیره است که مستقیماً از بازنمایی‌های نهفته اندیشمند، نشانه‌های
صوتی می‌سازد. این مدل برای پروژه ما به دو دلیل کنار گذاشته شد: نسخه هفت‌میلیارد‌پارامتری
آن به حافظه‌ای به‌مراتب بیش از بازه هدف \N{12} تا \N{24} گیگابایت نیاز
دارد، و فهرست‌های منتشرشده ورودی و خروجی گفتاری آن فارسی را شامل
نمی‌شوند.

\subsection{\en{Moshi} و انتخاب نهایی مسیر مستقیم}

از میان مدل‌های بررسی‌شده، \en{Moshi}~\cite{defossez2024moshi} تنها موردی بود
که هر سه شرط لازم برای این پروژه را برآورده می‌کرد: قابلیت تمام‌دوطرفه
«ذاتی» از طریق مدل‌سازی موازی دو جریان صوتی، وزن‌های متن‌باز، و
مهم‌تر از همه یک «آموزش‌دهنده رسمی تطبیق کم‌پارامتر» که به‌طور مستقیم
از سوی سازندگان مدل منتشر شده بود~\cite{moshifinetune2024}. همان‌طور که در
بخش~\ref{sec:codec} توضیح داده شد، سازوکار «تک‌گویی درونی» این مدل امکان
پایش جداگانه زیان متنی و زیان صوتی را نیز فراهم می‌کند که برای عیب‌یابی
آزمایش‌های ناموفق ارزش عملی داشت.

پیشینه نزدیک‌تر به این ایده،
\en{dGSLM}~\cite{nguyen2023dgslm} است که مدل‌سازی مولد گفت‌وگوی گفتاری را
با دو جریان موازی و بدون واسطه متنی معرفی کرد و نشان داد که پدیده‌هایی
مانند هم‌پوشانی و بازخورد کلامی را می‌توان به‌صورت آموخته‌شده بازتولید کرد.

\subsection{\en{PersonaPlex} و کنترل نقش در مدل تمام‌دوطرفه}

\en{PersonaPlex}~\cite{roy2026personaplex} یک مدل گفتار‌به‌گفتار تمام‌دوطرفه
بر پایه معماری \en{Moshi} است که با پرامپت ترکیبی متنی و نمونه صوتی، نقش و
صدای دستیار را شرطی می‌کند. وزن‌های متن‌باز آن منتشر شده‌اند و قابلیت
تمام‌دوطرفه ذاتی را حفظ می‌کنند. برای این پروژه اما دو محدودیت تعیین‌کننده
داشت: انتشار آن انگلیسی‌محور است و فارسی در آن پشتیبانی نشده است، و
سخت‌افزار آزمون‌شده مستند آن خارج از بازه \N{12} تا \N{24} گیگابایت
نیازمندی ن۶ است. همچنین یک آموزش‌دهنده رسمی برای تطبیق آن با زبانی جدید،
در زمان تصمیم‌گیری این پروژه، در دسترس نبود. بنابراین به‌عنوان گزینه
مقایسه در نظر گرفته شد و برای مسیر مستقیم انتخاب نشد.

\section{منابع گفتاری زبان فارسی}
\label{sec:persian-resources}

برای فارسی، وضعیت منابع در سه حوزه نامتقارن است.

در \مهم{بازشناسی گفتار}، مدل‌های قابل استفاده وجود دارند. پیکره چند‌زبانه
\en{Common Voice}~\cite{ardila2020commonvoice} بخش فارسی دارد، و مدل‌های
بومی‌سازی‌شده بر پایه معماری \en{Conformer} در بستر
\en{NeMo}~\cite{kuchaiev2019nemo} به کیفیت قابل قبولی رسیده‌اند. سرویس
بازشناسی به‌کاررفته در این پروژه از همین دسته است.

در \مهم{تبدیل متن به گفتار}، شاخص‌ترین منبع عمومی
\en{ManaTTS}~\cite{qharabagh2025manatts} است: بزرگ‌ترین پیکره عمومی
تک‌گوینده فارسی با حدود \N{86} ساعت صوت با بسامد نمونه‌برداری
\N{44.1} کیلوهرتز و مجوز آزاد، همراه با یک زنجیره پردازش متن‌باز شامل
ابزار بخش‌بندی جمله و یک روش هم‌ترازی اجباری طراحی‌شده برای زبان‌های
کم‌منبع. مدل صوتی فارسی مورد استفاده در این پروژه از همین منبع مشتق شده
است.

در \مهم{داده مکالمه‌ای تمام‌دوطرفه} اما شکاف جدی است. آنچه یک مدل
تمام‌دوطرفه برای یادگیری نیاز دارد، گفتار دو‌طرفه طبیعی همراه با
هم‌پوشانی، وقفه و بازخورد کلامی است. منابع فارسی موجود، عمدتاً تک‌گوینده و
خوانشی‌اند و این پدیده‌ها را در بر ندارند. پیکره‌های تلفنی گفت‌وگویی فارسی
در آرشیوهای پژوهشی موجودند اما مجوز محدود دارند و مقیاس آن‌ها کمتر از بازه
\N{100} تا \N{200} ساعت مورد نظر تعریف پروژه است. همین شکاف است که ساخت
مجموعه داده مشتق‌شده در فصل~\ref{ch:method} را ضروری کرد.

\section{تشخیص وقفه و نوبت‌گیری}
\label{sec:bargein-lit}

پیشینه این حوزه دو شاخه دارد.

شاخه \مهم{کلاسیک} بر ویژگی‌های آوایی دست‌ساز و آشکارسازهای آماری استوار
است. آشکارساز فعالیت گفتاری مبتنی بر مدل آماری~\cite{sohn1999vad} نمونه
مرجع این شاخه است. تشخیص وقفه در سامانه‌های تلفنی گفتاری سنتی نیز عملاً یک
آشکارساز فعالیت گفتاری با آستانه محافظه‌کارانه بود؛ مشکل شناخته‌شده این
سامانه‌ها آن بود که صدای بلندگو یا نویز محیط به‌اشتباه به‌عنوان وقفه
شناخته می‌شد. مرور جامع سازوکارهای نوبت‌گیری در سامانه‌های گفت‌وگو و تعامل
انسان و ربات در~\cite{skantze2021turntaking} آمده است، که نشان می‌دهد
نوبت‌گیری صرفاً یک مسئله تشخیص سکوت نیست و به نشانه‌های نوایی و نحوی
وابسته است. مبنای زبان‌شناختی این موضوع به کار کلاسیک~\cite{sacks1974turntaking}
بازمی‌گردد که نشان داد فاصله میان نوبت‌ها در گفت‌وگوی انسانی به‌طور شگفت‌آوری
کوتاه و منظم است.

شاخه \مهم{یادگیری‌شده} نوبت‌گیری را در خود مدل مولد جای می‌دهد. در
\en{Moshi} و \en{dGSLM} هیچ آشکارساز جداگانه‌ای وجود ندارد و نوبت‌گیری بخشی
از توزیع آموخته‌شده روی دو جریان موازی است. در \en{Freeze-Omni} یک راه
میانی به کار رفته است: یک سر طبقه‌بندی سبک روی مدل زبانی، حالت گفت‌وگو را
پیش‌بینی می‌کند.

تعریف پروژه صریحاً شاخه کلاسیک را خواسته است: یک ماژول بر پایه انرژی،
زیر‌و‌بمی و \en{MFCC} با هدف دقت بالای \N{80} درصد. پیاده‌سازی این پروژه
همین مسیر را دنبال می‌کند، با این تفاوت روشمند که به‌جای آستانه‌گذاری
دستی، یک طبقه‌بند تقویت گرادیانی روی بردار ویژگی پنجره متن آموزش می‌بیند و
آستانه تصمیم آن تنها بر نشست‌های اعتبارسنجی انتخاب می‌شود.

برای برچسب‌گذاری خودکار رخدادهای تعاملی در داده ضبط‌شده، از
\مهم{تفکیک گوینده} استفاده شده است. ابزارهای مرجع این حوزه
\en{pyannote.audio}~\cite{bredin2020pyannote} و روش‌های تفکیک چند‌مقیاسی با
وزن‌دهی پویا~\cite{park2022msdd} هستند؛ سرویس تفکیک گوینده مورد استفاده در
این پروژه از خانواده دوم است.

\section{ارزیابی خودکار کیفیت گفت‌وگو}
\label{sec:eval-lit}

سنجش کیفیت پاسخ گفت‌وگو با سنجه‌های تطابق رشته‌ای معنادار نیست. رویکردی که
در سال‌های اخیر مرجع شده است، به‌کارگیری یک مدل زبانی بزرگ به‌عنوان داور
است. کار مرجع این حوزه~\cite{zheng2023mtbench} نشان داد که داوری خودکار با
مدل‌های قوی، همبستگی بالایی با ترجیح انسانی دارد، اما در همان کار
اریبی‌های سامانه‌ای آن نیز مستند شده است: اریبی موقعیت، اریبی طول پاسخ و
اریبی خودپسندی نسبت به خروجی مدل‌های هم‌خانواده.

این پروژه از همین روش به‌عنوان «تنها» سنجه معنایی در دسترس استفاده
کرده است، اما با سه محدودسازی: دوسوکورسازی و تصادفی‌سازی ترتیب بازوها،
کدگشایی قطعی، و تکرار هر فراخوانی داور با گزارش نرخ توافق. مهم‌تر آن‌که
چون داور و نامزد در این پروژه از یک خانواده مدل‌اند، نتیجه صریحاً به‌عنوان
یک «تقریب خودکار هم‌خانواده» برچسب‌گذاری شده و جانشین ارزیابی انسانی یا
یک معیار مستقل شمرده نشده است. برای سنجش انسانی، پروتکل مرجع، امتیاز
میانگین نظر بر پایه توصیه‌نامه \en{ITU-T P.800}~\cite{itu1996p800} است که در
این پروژه طراحی شد اما اجرا نشد.

\section{جایگاه این پروژه}
\label{sec:positioning}

جدول~\ref{tab:positioning} جایگاه این پروژه را در برابر کارهای پیشین
خلاصه می‌کند.

\begin{table}[ht]
\centering
\caption{مقایسه این پروژه با کارهای پیشین از نظر قابلیت‌های مورد نیاز
         تعریف پروژه}
\label{tab:positioning}
\small
\begin{tabular}{L{2.5cm}ccccc}
\toprule
\textbf{کار} & \textbf{تمام‌} & \textbf{وزن} & \textbf{آموزش‌دهنده} &
\textbf{پشتیبانی} & \textbf{بازه} \\
             & \textbf{دوطرفه} & \textbf{متن‌باز} & \textbf{رسمی} &
\textbf{فارسی} & \textbf{حافظه} \\
\midrule
\en{Mini-Omni2}   & فرمان‌محور & دارد & ندارد & ندارد & مناسب \\
\en{LLaMA-Omni2}  & ندارد & دارد & ندارد & ندارد & مناسب \\
\en{Freeze-Omni}  & دارد & دارد & جزئی & ندارد & مناسب \\
\en{Qwen2.5-Omni} & دارد & دارد & ندارد & ندارد & نامناسب \\
\en{Moshi}        & ذاتی & دارد & دارد & ندارد & نامناسب \\
\en{PersonaPlex}  & ذاتی & دارد & جزئی & ندارد & نامناسب \\
\midrule
این پروژه، مسیر آبشاری & کنترل بیرونی & دارد & بی‌نیاز & دارد & مناسب \\
این پروژه، مسیر مستقیم & ذاتی & دارد & دارد & ناموفق & نامناسب \\
\bottomrule
\end{tabular}
\end{table}

\مرزادعا{ستون «بازه حافظه» به تناسب با بازه \N{12} تا \N{24} گیگابایت مورد
نظر تعریف پروژه اشاره دارد و برای مسیر مستقیم این پروژه بر پایه حافظه
اندازه‌گیری‌شده آموزش روی سخت‌افزار موجود گزارش شده است، نه یک اندازه‌گیری
روی سخت‌افزار هدف.}

از این مقایسه سه نکته برمی‌آید.

نخست، هیچ‌یک از مدل‌های سرتاسری موجود از فارسی پشتیبانی نمی‌کند؛ بنابراین
هر مسیر مستقیمی ناچار از تطبیق است و تطبیق نیز به آموزش‌دهنده رسمی نیاز
دارد. تنها \en{Moshi} این شرط را به‌طور کامل برآورده می‌کرد.
\en{PersonaPlex} نیز تمام‌دوطرفه ذاتی دارد اما انگلیسی‌محور است و آموزش‌دهنده
تطبیق آن برای فارسی در دسترس نبود.

دوم، مدل‌هایی که قابلیت تمام‌دوطرفه ذاتی دارند، در بازه حافظه هدف جای
نمی‌گیرند. این تعارض میان نیازمندی ن۲ و ن۶ در تعریف پروژه، یک محدودیت
ساختاری است و نه یک نقص پیاده‌سازی.

سوم، سهم متمایز این پروژه در جای دیگری است. کارهای پیشین، «مدل»
منتشر می‌کنند؛ این پروژه، برای یک زبان کم‌منبع مشخص، «یک زنجیره
داده و یک چارچوب ارزیابی شواهد‌محور» می‌سازد و نشان می‌دهد کدام ادعاها با
این شواهد قابل اثبات و کدام‌ها غیرقابل اثبات‌اند. به‌طور خاص، نتیجه منفی
مسیر مستقیم، که در بخش~\ref{sec:res-moshi} گزارش شده است، یک یافته
قابل استناد برای پژوهشگران بعدی است: کاهش زیان تحت تحمیل معلم، در تطبیق
کم‌پارامتر یک مدل گفتار‌به‌گفتار به زبانی جدید، معیار کافی برای انتخاب
چک‌پوینت نیست.
